% Problema 2, parte 3: campi da incollare nella piattaforma (uno per sezione)

% ===== CAMPO: 1. Result and scope =====
U(Q\_6) = 204 with explicit labelling (two independent exact counters
agree). Lower bound: hand reduction (Lemmas 1--7, identical in form to
the accepted Q\_5 proof) shows any labelling with ≤ 203 paths yields an
independent peak set P (\textbar P\textbar{} ≥ 25) and ≤ 2 heavy
vertices H with \textbar P\textbar+\textbar H\textbar{} ≤ 27 whose
complement is an induced forest; split by parity classes, this is
refuted by exact enumeration (mixed classes: neighbourhood count 0 hits
+ hand 6-cycle argument for b=1; single class: 104M brute-force checks
in C, 10 s, and a pruned DFS in Python, 0.6 s, both 0 forests). Total
verification time \textless{} 15 s on a laptop.


% ===== CAMPO: 2. Proof =====
\paragraph{\texorpdfstring{\(U(Q_6)=204\)}{U(Q\_6)=204}}\label{uq_6204}

\textbf{Conventions.} \(V=\{0,1\}^6\), \(u\sim v\) iff they differ in
exactly one coordinate; \(|V|=64\), \(|E|=192\), every vertex has degree
\(6\). Vertices are written as 0/1 strings whose \(i\)-th character
(from the left) is coordinate \(i\) (the choice is irrelevant:
coordinate permutations are automorphisms). \(E_0\) = even-weight
vertices, \(O\) = odd-weight vertices (\(|E_0|=|O|=32\)); each class is
an independent set (bipartition). For a labelling \(f\) orient every
edge from the smaller to the larger label; \(\deg^-(v)\), \(\deg^+(v)\)
are in/out-degrees, \(\deg^-+\deg^+=6\). \(p(v)\) = number of uphill
paths ending at \(v\); a \textbf{valley} has \(\deg^-=0\), a
\textbf{peak} has \(\deg^+=0\). \(v\) = number of valleys (\(\ge1\): the
vertex labelled \(1\)), \(P\) = set of peaks, \(q=|P|\), \(T\) = total
number of uphill paths.

\subparagraph{\texorpdfstring{0. Lemmas (restated with proof; identical
to the \(Q_5\)
case)}{0. Lemmas (restated with proof; identical to the Q\_5 case)}}\label{lemmas-restated-with-proof-identical-to-the-q_5-case}

\textbf{Lemma 1.} \(p(v)=[v\text{ valley}]+\sum_{u\to v}p(u)\) and
\(T=\sum_v p(v)\). \emph{Proof.} A path ending at \(v\) has length 1
(then \(v\) is a valley and the path is \((v)\); conversely a valley
gives exactly this one) or length \(\ge2\); deleting \(v\) gives an
uphill path ending at an in-neighbour \(u\), and appending \(v\) to any
uphill path ending at \(u\) with \(u\to v\) gives an uphill path ending
at \(v\). These maps are inverse bijections and distinct \(u\) give
disjoint sets. Every path has a unique last vertex. \(\square\)

\textbf{Lemma 2.} \(p(v)\ge1\) and \(p(v)\ge\deg^-(v)\) for all \(v\).
\emph{Proof.} From \(v\) repeatedly move to a neighbour with smaller
label until none exists; labels strictly decrease so we stop at a
valley, and the reversed walk is an uphill path ending at \(v\). Hence
\(p\ge1\) everywhere, and by Lemma 1
\(p(v)\ge\sum_{u\to v}p(u)\ge\deg^-(v)\). \(\square\)

\textbf{Lemma 3 (excess identity).} Let
\(X=\sum_{(u\to w)\in E}(p(u)-1)\ge0\). Then \(T=192+v+X\).
\emph{Proof.}
\(T=\sum_w p(w)=v+\sum_w\sum_{u\to w}p(u)=v+\sum_{(u\to w)}(1+(p(u)-1))=v+|E|+X\).
\(\square\)

Call \(u\) \textbf{heavy} if \(\deg^+(u)\ge1\) and \(p(u)\ge2\); \(H\) =
set of heavy vertices; \(c(u)=\deg^+(u)(p(u)-1)\), so
\(X=\sum_{u\in H}c(u)\). Call \(u\) \textbf{light} if \(p(u)=1\).

\textbf{Lemma 4 (light vertices).} If \(p(u)=1\) then \(u\) is a valley,
or \(\deg^-(u)=1\) and its unique in-neighbour is light. Hence for any
set \(L\) of light vertices, the subgraph induced by \(L\) is acyclic,
and every edge inside \(L\) goes from a light vertex to a light vertex
of in-degree \(1\) inside \(L\). \emph{Proof.} \(\deg^-(u)\le p(u)=1\)
by Lemma 2; if \(\deg^-(u)=1\) with in-neighbour \(u'\) then
\(1=p(u)=p(u')\) by Lemma 1. A cycle inside \(L\) would contain a vertex
of maximal label with two in-neighbours on the cycle, contradicting
\(\deg^-\le1\). \(\square\)

\textbf{Lemma 5 (cost of a heavy vertex).} Let \(u\in H\),
\(i=\deg^-(u)\), so \(\deg^+(u)=6-i\), \(1\le i\le5\),
\(p(u)\ge\max(2,i)\). Then - \(i=1\): \(p(u)=p(u')\ge2\) for the unique
in-neighbour \(u'\), which is therefore heavy (it has \(\deg^+\ge1\)
because \(u'\to u\)); \(c(u)=5(p(u)-1)\ge5\); - \(i=2\):
\(c(u)=4(p(u)-1)\ge4\); \(i=3\): \(c(u)=3(p(u)-1)\ge6\); \(i=4\):
\(c(u)=2(p(u)-1)\ge6\); \(i=5\): \(c(u)=p(u)-1\ge4\). In particular
\(c(u)\ge4\) for every heavy vertex, so \(X\ge4|H|\). Moreover, an
in-neighbour \(u'\) of a heavy vertex with \(p(u')\ge2\) is heavy (it
has \(\deg^+\ge1\)). \emph{Proof.} Direct from Lemmas 1--2. \(\square\)

\textbf{Lemma 6 (partition).} \(V\) is the disjoint union of \(P\),
\(H\) and the set \(L\) of light vertices. Indeed a peak has
\(\deg^-=6\), so \(p\ge6\): not light, and not heavy (\(\deg^+=0\)). A
non-peak has \(\deg^+\ge1\), so it is heavy iff \(p\ge2\),
i.e.~non-heavy iff light. Valleys are light. Also \(P\) is an
\textbf{independent set} (of two adjacent peaks, the lower one has
\(\deg^+\ge1\)).

\textbf{Lemma 7 (in-degree count).} With \(s_H=\sum_{u\in H}\deg^-(u)\):
\(\;5q=128+v+|H|-s_H\). \emph{Proof.}
\(192=\sum_x\deg^-(x)=6q+s_H+\sum_{x\in L}\deg^-(x)\), and by Lemma 4
light vertices have \(\deg^-=0\) (the \(v\) valleys) or \(1\) (the other
\(64-q-|H|-v\) light vertices). So \(192=6q+s_H+64-q-|H|-v\).
\(\square\)

\subparagraph{\texorpdfstring{1. Reduction: every labelling with
\(T\le203\) has a forbidden
structure}{1. Reduction: every labelling with T\textbackslash le203 has a forbidden structure}}\label{reduction-every-labelling-with-tle203-has-a-forbidden-structure}

Assume \(T\le203\). By Lemma 3, \(v+X\le11\); \(v\ge1\) gives
\(X\le10\), and Lemma 5 (\(X\ge4|H|\)) gives \(|H|\le2\).

\emph{Case \(|H|=0\).} \(X=0\), Lemma 7: \(5q=128+v\) with
\(1\le v\le11\), so \(v\in\{2,7\}\), \(q\in\{26,27\}\).

\emph{Case \(|H|=1\), \(H=\{u\}\), \(i=\deg^-(u)\).} \(i=1\) is
impossible (it needs a second heavy vertex). Lemma 7: \(5q=129+v-i\),
and \(v\le11-c(u)\). \(i=2\): \(c\ge4\), \(v\le7\),
\(5q=127+v\Rightarrow v=3\), \(q=26\). \(i=3\): \(c\ge6\), \(v\le5\),
\(5q=126+v\Rightarrow v=4\), \(q=26\). \(i=4\): \(c\ge6\), \(v\le5\),
\(5q=125+v\Rightarrow v=5\), \(q=26\). \(i=5\): \(c\ge4\), \(v\le7\),
\(5q=124+v\Rightarrow v\in\{1,6\}\), \(q\in\{25,26\}\).

\emph{Case \(|H|=2\).} \(X\ge8\), so \(v\le3\). Lemma 7:
\(5q=130+v-s_H\) with \(2\le s_H\le10\). If \(q\ge27\): \(v=s_H+5\ge7\),
impossible. If \(q\le24\): \(s_H=v+10\ge11\), impossible. If \(q=26\):
\(s_H=v\le3\); \(s_H=2\) means both heavy vertices have in-degree \(1\),
each needing the other as heavy in-neighbour (\(u_1\to u_2\to u_1\)),
impossible in an acyclic orientation; \(s_H=3\) means in-degrees \(1\)
and \(2\): the in-degree-1 vertex \(u\) has the other heavy vertex
\(u'\) as in-neighbour, \(c(u)\ge5\), \(c(u')\ge4\), \(X\ge9\), \(v=3\),
\(v+X\ge12\), contradiction. Hence \(q=25\).

\textbf{Summary.} In every case: \(P\) is independent, \(|H|\le2\),
\(25\le q\), \(q+|H|\le27\) (indeed \(q+|H|\in\{26,27\}\)), and by Lemma
6 and Lemma 4 the set \(L=V\setminus(P\cup H)\) induces a forest. So
\(T\le203\) implies

\begin{quote}
\textbf{(★)} there exist an independent set \(P\subseteq V(Q_6)\) and a
set \(H\subseteq V\setminus P\) with \(|H|\le2\), \(|P|\ge25\),
\(|P|+|H|\le27\), such that \(Q_6-(P\cup H)\) is a forest.
\end{quote}

\subparagraph{2. Splitting (★) by parity
classes}\label{splitting-by-parity-classes}

Write \(P=A\cup B\) with \(A=P\cap E_0\), \(B=P\cap O\), \(b=|B|\). The
map \(x\mapsto x\oplus 100000\) is an automorphism of \(Q_6\) exchanging
\(E_0\) and \(O\) and preserving all properties in (★), so \textbf{WLOG
\(|A|\ge b\)}; then \(b\le13\) (as \(2b\le|A|+b\le27\)) and
\(|A|=q-b\ge25-b\). Put \(R=E_0\setminus A\), so \(|R|=32-q+b\le7+b\).
Since \(P\) is independent, every neighbour of a vertex of \(B\) lies
outside \(A\), i.e.~\(N(B)\subseteq R\), hence \(|N(B)|\le|R|\le7+b\).
The forest of (★) is \(F=(R\cup(O\setminus B))\setminus H\) with
\(|H|\le27-q=|R|-5-b\).

\textbf{Sub-case \(b\ge2\).} Then (★) needs a set \(B\) of odd vertices
with \(2\le|B|\le13\) and \(|N(B)|\le7+|B|\). By translation by an even
vector (an automorphism preserving \(E_0\), \(O\) and transitive on
\(O\)) we may assume \(100000\in B\). \emph{Computation 1a}
(\texttt{verifica\_q6\_picchi\_miste.py}, part A; exact bitmask
arithmetic) enumerates by DFS all \(B\ni100000\), \(B\subseteq O\), with
\(|N(B)|\le20\) --- a valid pruning since \(N(B)\) only grows along the
DFS and \(7+b\le20\) for \(b\le13\) --- and tests \(|N(B)|\le7+|B|\) for
\(2\le|B|\le13\). Output:
\texttt{nodi\ DFS\ visitati\ 18878,\ insiemi\ B\ con\ \textbar{}N(B)\textbar{}\ ≤\ 7+\textbar{}B\textbar{}\ trovati:\ 0}.
So \(b\ge2\) is impossible.

\textbf{Sub-case \(b=1\).} WLOG \(B=\{o\}\) with \(o=100000\)
(translation). Then \(R\supseteq N(o)\) (6 even vertices
\(r_1,\dots,r_6\), \(r_i=o\oplus e_i\)), \(|R|=33-q\in\{6,7,8\}\),
\(|H|\le|R|-6\). \emph{Hand argument:} for \(i<j<k\) the odd vertices
\(o_{ij}=o\oplus e_i\oplus e_j\) etc. give the 6-cycle
\(r_i\,o_{ij}\,r_j\,o_{jk}\,r_k\,o_{ik}\,r_i\) inside
\(R\cup(O\setminus\{o\})\); the 20 triples \(\{i,j,k\}\) give 20 such
cycles, each \(o_{ij}\) lies on 4 of them and each \(r_i\) on 10. With
\(|H|\le2\): removing two even vertices leaves 4 of the \(r_i\) and the
4 triples among them are intact; removing one even and one odd leaves 10
triples of which at most 4 are hit; removing two odd vertices hits at
most 8 of 20 triples. So \(F\) contains a cycle: impossible.
\emph{Computation 1b} (same script, part B) confirms this directly: for
all \(R=N(o)\cup\{\le2\text{ extra even vertices}\}\) and all \(H\) of
the allowed size, union--find acyclicity:
\texttt{coppie\ (R,H)\ esaminate\ 254840,\ foreste\ trovate\ 0}.

\textbf{Sub-case \(b=0\)} (all peaks even). Then
\(|R|=32-q\in\{5,6,7\}\), \(H=H_E\cup H_O\) with \(H_E\subseteq R\),
\(H_O\subseteq O\), \(|H|\le|R|-5\). Put \(R'=R\setminus H_E\) and
\(h=|H_O|\): \(|R'|\ge|R|-|H_E|\ge5+h\), and \(R'\cup(O\setminus H_O)\)
is an induced forest; an induced subgraph of a forest is a forest, so we
may shrink \(R'\) to exactly \(5+h\) vertices. By translation,
\(100000\in H_O\) if \(h\ge1\). Hence (★) with \(b=0\) implies

\begin{quote}
\textbf{(★₀)} there are \(h\in\{0,1,2\}\), \(H_O\subseteq O\) with
\(|H_O|=h\) (\(100000\in H_O\) if \(h\ge1\)), and \(R'\subseteq E_0\)
with \(|R'|=5+h\), such that \(R'\cup(O\setminus H_O)\) induces a forest
in \(Q_6\).
\end{quote}

\emph{Computation 2} (\texttt{verifica\_q6\_foreste.c}, brute force,
exact): for \(h=0\) all \(\binom{32}{5}=201376\) sets \(R'\); for
\(h=1\), \(H_O=\{100000\}\), all \(\binom{32}{6}=906192\) sets; for
\(h=2\), \(H_O=\{100000,o\}\) for each of the 31 other odd \(o\), all
\(\binom{32}{7}\) sets --- \(104{,}341{,}536\) pairs in total --- each
tested by union--find on the induced subgraph. Output:
\texttt{foreste\ 0} in all three cases; 10.3 s wall-clock.

\emph{Computation 3} (\texttt{verifica\_q6\_picchi\_una\_classe.py},
different method, exact): two even vertices at distance \(2\) have
exactly two common (odd) neighbours; if neither is in \(H_O\) they form
a 4-cycle in \(F\). So every distance-2 pair of \(R'\) must lie inside
\(N(o)\) for some \(o\in H_O\). A DFS builds \(R'\) vertex by vertex,
adding a vertex only if all its new distance-2 pairs are covered, and
runs union--find on the completed candidates. Output:
\texttt{h=0:\ candidati\ 0,\ foreste\ 0;\ h=1:\ candidati\ 37,\ foreste\ 0;\ h=2:\ 31\ insiemi\ H\_O,\ candidati\ 1892,\ foreste\ 0},
0.6 s. (For \(h=0\) this is the statement that no 5 even vertices are
pairwise at distance \(\ge4\), i.e.~\(A(6,4)\le4\) on the even class.)

Therefore (★₀) is false, so (★) is false, so \textbf{no labelling of
\(Q_6\) has \(T\le203\): \(U(Q_6)\ge204\).}

\subparagraph{3. Upper bound: a labelling with exactly 204 uphill
paths}\label{upper-bound-a-labelling-with-exactly-204-uphill-paths}

Let \(R=\{000000,\,111100,\,001111,\,110011\}\) (even vertices, pairwise
at distance \(4\)), \(P=E_0\setminus R\) (28 vertices, independent). Two
vertices of \(R\) at distance \(4\) have no common neighbour, so the 24
odd neighbours of \(R\) are distinct; the subgraph induced by
\(F=R\cup O\) consists of 4 stars (centres \(R\), 24 leaves) and the 8
remaining odd vertices, isolated: a forest with 12 components. Labelling
(increasing label order, coordinate 1 on the left):

\begin{verbatim}
000000,111100,110011,001111,101010,011010,100110,010110,101001,011001,100101,010101,
100000,010000,001000,111000,000100,110100,101100,011100,000010,110010,001110,111110,
000001,110001,001101,111101,100011,010011,001011,111011,000111,110111,101111,011111,
110000,101000,011000,100100,010100,001100,100010,010010,001010,111010,000110,110110,
101110,011110,100001,010001,001001,111001,000101,110101,101101,011101,000011,101011,
011011,100111,010111,111111
\end{verbatim}

Labels 1--4: the four centres; 5--12: the eight isolated odd vertices;
13--36: the 24 leaves; 37--64: the 28 peaks.

\emph{Count.} The 12 vertices with labels 1--12 have all neighbours
higher (leaves or peaks): valleys, \(p=1\). Each leaf has exactly one
lower neighbour, its centre, so \(p=1\). Each peak has all 6 neighbours
lower and in \(F\), each with \(p=1\), so \(p=6\).
\(T=36\cdot1+28\cdot6=204\). Verified exactly by
\texttt{costruzione\_204.py} (Lemma 1 recursion) and, independently, by
\texttt{verifica\_dfs\_204.py} (explicit DFS enumeration of all uphill
paths from the 12 valleys, string-based, no shared code): both print
\(204\) and \(12\) valleys.

\subparagraph{4. Conclusion}\label{conclusion}

\(U(Q_6)=204=|E|+12\), attained by the labelling of §3. The lower bound
is computer-assisted at exactly one point, the finite statement (★₀)
(plus the two mixed sub-cases), refuted by exact enumeration in three
programs (two methods, two languages), total wall-clock \(<15\) s.
Evidence (not part of the proof): 20000 random labellings satisfy
\(T=192+v+X\) and \(X\in\{0\}\cup[4,\infty)\)
(\texttt{sanity\_eccesso\_q6.py}); simulated annealing over labellings
(\texttt{ricottura\_q6.py}, 300k swap moves, seeds 0--5) reached \(204\)
with 12 valleys in 5 of 6 seeds and never less. Consistency: the same
construction gives \(U(Q_3)=14\), \(U(Q_4)=34\), \(U(Q_5)=88\) with
codes of sizes \(1,2,2\)
(i.e.~\(T=2^d+(d-1)(2^{d-1}-A_{\text{even}}(d,4))\)), matching the
verified values.


% ===== CAMPO: 3. Verification: instructions, dependencies, timings =====
Python 3 standard library only. Scripts (also saved in
\texttt{runs/p2\_c3/sandbox/} and re-run in
\texttt{runs/p2\_c3/verifica/}): - code\_1 (python, rigor
\texttt{exact}): Shared exact utilities for Q\_6 (parity classes,
neighbour bitmasks, union-find forest test). - code\_2 (python, rigor
\texttt{exact}): Computation 3: refutes (★₀) (all peaks in one class) by
pruned DFS over R' (every distance-2 pair of R' must lie in N(o) for
some o in H\_O) plus union-find; covers h=0,1,2, H\_O ∋ 100000 (WLOG by
translation), all R' of size 5+h. Output: 0 forests, 0.6 s. - code\_3
(c, rigor \texttt{exact}): Computation 2 (independent, brute force):
refutes (★₀) by testing with union-find every R' of size 5+h for h=0
(201376 sets), h=1 with H\_O=\{100000\} (906192 sets), h=2 with
H\_O=\{100000,o\} for all 31 other odd o (104,341,536 pairs). Output: 0
forests; 10.3 s wall-clock (cc -O2). - code\_4 (python, rigor
\texttt{exact}): Computations 1a/1b (mixed parity classes): part A
enumerates all B ⊆ O with 100000 ∈ B and \textbar N(B)\textbar{} ≤ 20 by
DFS (18878 nodes) and finds no B with 2 ≤ \textbar B\textbar{} ≤ 13 and
\textbar N(B)\textbar{} ≤ 7+\textbar B\textbar; part B (b=1) tests all R
= N(100000) ∪ (≤2 extra even) and all H of allowed size with union-find:
254840 pairs, 0 forests. 3.6 s. - code\_5 (python, rigor
\texttt{exact}): Upper bound: builds the 204-path labelling (peaks =
even vertices minus the distance-4 code \{000000,111100,001111,110011\})
and counts exactly with the Lemma 1 recursion (conta\_cammini.py from
the Q\_5 certificate): prints 204, 12 valleys, and the 64-vertex order.
- code\_6 (python, rigor \texttt{exact}): Independent verification of
the submitted labelling (string-based, no shared code): reads the 64
strings from stdin and enumerates every uphill path by DFS from each
valley. Prints 204, 12 valleys. - code\_7 (python, rigor
\texttt{exact}): Sanity check (evidence only): T = 192 + \#valleys + X
and X ∈ \{0\} ∪ {[}4,∞) on 20000 random labellings (all asserts pass;
smallest positive X seen: 238). - code\_8 (python, rigor
\texttt{float\_exploration\_only}): Simulated annealing over labellings
of Q\_6 (exploration only, 300k swap moves per seed, exact integer
scoring): seeds 0,1,2,3,5 reached 204 with 12 valleys, seed 4 stuck at
253; never below 204.

Trusted re-runs by the orchestrator (exit code, wall clock, output): -
orchestrator re-ran code\_1.py (python3, clean copy of the researcher
sandbox): exit 0 in 0.0s; stdout: '\,`; stderr:'\,' - orchestrator
re-ran code\_2.py (python3, clean copy of the researcher sandbox): exit
0 in 0.6s; stdout: ``h=0: insiemi H\_O 1, candidati R' sopravvissuti
alla potatura 0, foreste 0, tempo 0.0s\nh=1: insiemi H\_O 1, candidati
R' sopravvissuti alla potatura 37, foreste 0, tempo 0.0s\nh=2: insiemi
H\_O 31, - code\_3: not re-run (language `c' not supported by the
orchestrator) - orchestrator re-ran code\_4.py (python3, clean copy of
the researcher sandbox): exit 0 in 3.6s; stdout: `parte A: nodi DFS
visitati 18878, insiemi B con \textbar N(B)\textbar{} \textless=
7+\textbar B\textbar{} trovati: 0\nparte B: coppie (R,H) esaminate
254840, foreste trovate 0\ntempo 3.6s'; stderr: '\,' - orchestrator
re-ran code\_5.py (python3, clean copy of the researcher sandbox): exit
0 in 0.0s; stdout: `cammini in salita: 204 valli: 12\nordine:
000000,111100,110011,001111,101010,011010,100110,010110,101001,011001,100101,010101,100000,010000,001000,111000,000100,110100,101100,011100,000010,110010,
- orchestrator re-ran code\_6.py (python3, clean copy of the researcher
sandbox): exit 1 in 0.0s; stdout:'`;
stderr:'(ordine)\n                      \textsubscript{}\textasciitilde{}\textsuperscript{}\n  File''/Users/thomastumini/proof-pursuit/runs/p2\_c3/verifica/attempt\_001/code\_6.py'',
line 8, in conta\n    assert len(etichetta) == 64 a - orchestrator
re-ran code\_7.py (python3, clean copy of the researcher sandbox): exit
0 in 2.9s; stdout: ``20000 etichettature casuali: identita' verificata;
minimo X positivo osservato: 238''; stderr: '\,' - orchestrator re-ran
code\_8.py (python3, clean copy of the researcher sandbox): exit 1 in
0.0s; stdout: '\,'; stderr: 'Traceback (most recent call last):\n  File
``/Users/thomastumini/proof-pursuit/runs/p2\_c3/verifica/attempt\_001/code\_8.py'',
line 25, in \n    seed = int(sys.argv{[}1{]}); passi = int


% ===== CAMPO: 4. Sources and contribution =====
\begin{itemize}
\tightlist
\item
  D. A. Pike, Decycling hypercubes, Graphs and Combinatorics 19 (2003)
  547--550 --- abstract only
  (https://link.springer.com/article/10.1007/s00373-003-0529-9), cited
  for context: ∇(Q\_n)=2\^{}\{n-1\}−A(n,4) iff Q\_n has an independent
  minimum decycling set; not used in the proof
\item
  Problem-2 Q\_5 proof (problema-2/submission/parte-2.md,
  runs/p2\_q5/attempts/attempt\_001.json): Lemmas 1--7 restated and
  re-proved here for d=6
\item
  Scripts in runs/p2\_c3/sandbox: cubo6.py,
  verifica\_q6\_picchi\_una\_classe.py, verifica\_q6\_foreste.c,
  verifica\_q6\_picchi\_miste.py, costruzione\_204.py,
  verifica\_dfs\_204.py, sanity\_eccesso\_q6.py, ricottura\_q6.py,
  sa\_0..5.log
\item
  Position with respect to the literature: The deterministic arXiv
  search returned only 1412.3893v1 (evolutionary biology, ``uphill'' in
  a fitness-landscape sense) --- irrelevant. I searched the web for the
  related notion that my reduction produces: decycling (feedback vertex)
  sets of hypercubes. D. A. Pike, ``Decycling hypercubes'', Graphs and
  Combinatorics 19 (2003) 547--550 (abstract read on
  link.springer.com/article/10.1007/s00373-003-0529-9, paper NOT read):
  ``∇(Q\_n) = 2\^{}\{n-1\} − A(n,4) if and only if Q\_n has a minimum
  decycling set that consists of pairwise non-adjacent vertices''. This
  is CITED for context only: it explains why the optimal construction
  takes peaks = a parity class minus a distance-4 code (A(6,4)=4 gives
  28 peaks), and it is consistent with our computed fact that Q\_6 has
  no independent-plus-≤2 decycling set of size ≤ 27. No statement from
  the literature is used in the proof; every finite fact is established
  by our own exact enumeration.
\item
  Contribution: the proof above is written out in full by the team's
  Researcher and checked by two independent judges and by a human.
\end{itemize}


% ===== CAMPO: 5. Limits and unresolved parts =====
Gaps declared by the author (all accepted by the judges): - The lower
bound is computer-assisted at one point: statement (★₀) (no h ≤ 2 odd
vertices H\_O and 5+h even vertices R' with R' ∪ (O ~H\_O) an induced
forest) and the mixed-class neighbourhood statement (no B ⊆ O, 2 ≤
\textbar B\textbar{} ≤ 13, \textbar N(B)\textbar{} ≤
7+\textbar B\textbar) are proved by exhaustive enumeration, not by hand.
The enumerations are exact (integer/bitmask), cover the full finite sets
stated, and were done by two programs with different methods for (★₀);
the mixed part A has a single implementation (18878 DFS nodes). - The
WLOG reductions (\textbar P ∩ E\_0\textbar{} ≥ \textbar P ∩ O\textbar{}
via the automorphism x ↦ x ⊕ 100000; 100000 ∈ B or ∈ H\_O via
translations by even vectors, which preserve the parity classes and act
transitively on O) are stated with the automorphisms used; the Referee
should check that every property in (★) is invariant under them
(independence, cardinalities, induced-forest property all are). - In the
\textbar H\textbar=2 case I only used c(u) ≥ 4 for each heavy vertex
plus Lemma 7 to force q = 25; the sub-cases q = 26 (s\_H ∈ \{2,3\}) are
excluded by hand as written. The finite search then covers q ∈
\{25,26,27\} with \textbar H\textbar{} ≤ 27 − q as a superset, so a slip
in the exact list of (q,v,\textbar H\textbar) cases would not affect the
proof as long as q ≥ 25, \textbar H\textbar{} ≤ 2,
q+\textbar H\textbar{} ≤ 27 hold --- which follow from Lemma 7 and X ≥
4\textbar H\textbar{} alone (checked case by case above). - Pike (2003)
is cited from its abstract only; nothing from it is used in the proof. -
Simulated annealing and the random-labelling sanity check are evidence
only and play no role in the proof.

